\section{Related work}
\label{sec:related}

For each lineage we state what is inherited and where caging differs. No priority or novelty over supervisory control, runtime enforcement, or barrier methods is claimed.

\paragraph{Supervisory control of discrete-event systems.}
Ramadge and Wonham~\citep{ramadge1987,ramadge1989}, with the textbook of Cassandras and Lafortune~\citep{cassandras2008}: a plant generates a language of event sequences and a supervisor disables a subset of controllable events so that the closed-loop language stays within a specification. Almost the entire two-condition structure is inherited. Non-blocking is goal reachability; the safety specification is unsafe non-reachability; the supervisor is the cage; controllable versus uncontrollable events are edges a cage can remove versus edges it cannot; and the composition result of \S\ref{sec:composition} is a special case of the known fact that the conjunction of non-blocking supervisors need not be non-blocking, which is why modular supervisory control treats non-conflict as a separate property~\citep{wonham1988modular}. Caging differs in three respects that are the content of the framework: the plant is a learned policy whose transition relation is not given and is approximated by an abstraction; the state includes knowledge and memory, for which the classical theory has no counterpart; and the adjudicator is given its own requirements, since membership in $G$ is not observable without a verifier. Caging is supervisory control carried over to a cognitive architecture, with plant model, state, and observer each generalised.

\paragraph{Runtime enforcement, security automata, shielding.}
Schneider's execution monitors~\citep{schneider2000} enforce exactly the safety properties in the sense of Alpern and Schneider~\citep{alpern1985} and cannot enforce liveness; edit automata~\citep{ligatti2005} enlarge the class somewhat without crossing the safety/liveness boundary. Shield synthesis~\citep{bloem2015shield} and shielding for reinforcement learning~\citep{alshiekh2018} construct minimal-interference monitors from temporal specifications. The enforcement mechanism of any real cage is a monitor of this kind. Read through \S\ref{sec:conditions}, the classical result is the formal reason that guardrails accrete toward uselessness: a monitor cannot see the goal. Shields are additionally per-step, where a cage is a statement about the closure. Caging keeps the monitor as enforcement layer and adds the goal-side condition the monitor cannot express, checked at step 3 of \S\ref{sec:pipeline} before the monitor is deployed.

\paragraph{Control barrier functions and set invariance.}
A control barrier function renders a set forward-invariant by constraining the input~\citep{ames2017cbf,ames2019cbf}, with roots in Nagumo's theorem~\citep{nagumo1942} and viability theory~\citep{aubin1991}; control Lyapunov functions supply the goal-directed dual, and the combined quadratic program trades the two off at each instant. The pairing is inherited: a barrier alone is unsafe non-reachability with the one-sidedness of a security automaton, and the slack between barrier and descent direction is the continuous analogue of tightness. Here the state is discrete and partly symbolic, the dynamics are not differentiable, and the controller is a learned policy that is constrained rather than synthesised.

\paragraph{Barrier analysis in safety engineering.}
Haddon's energy-barrier model~\citep{haddon1973}, Reason's layered imperfect defences~\citep{reason1990}, defence in depth, and Leveson's STAMP~\citep{leveson2011}, which reframes safety as a control problem. The word cage is borrowed in spirit from this tradition, and so is its central lesson: barriers fail in combination, through interactions the designer of each barrier did not consider. Leveson's move from component reliability to control is the move of \S\ref{sec:problem} from the model to the system. Safety engineering has no goal-side condition (a plant that is safely shut down is a success); the goal condition is what distinguishes caging from a pure safety discipline.

\paragraph{Agent-side guardrails, sandboxes, permissions.}
Output and action guardrails, sandboxes, capability-based tool access, rule-based steering, and explicit success criteria are each an implementation of one cage: guardrails are Action Cages at the action level, sandboxes at the environment level, permission systems the held-permission component of the state, specification practices the Specification Cage. None of them, individually or as a list, states a goal condition or has a notion of tightness. The framework adds the explanation of why such a list tends to produce a safe and useless agent (\S\ref{sec:conditions}) and the classification of Table~\ref{tab:techniques}.

\paragraph{Memory governance and memory security.}
The works named are those examined, not a survey. Memory systems~\citep{memgpt,mem0,zep,amem,memos,evermemos} are retrieval-centred layers; Zep's bitemporal invalidate-do-not-delete rule is a precedent for traceability at the memory level. Memory attacks~\citep{agentpoison,minja} break the assumption that whatever is written to memory is true, the assumption a Memory Cage rejects. Governance layers~\citep{memlineage,collabmem,mapleguard} attach provenance and enforce access policy at write time; detectors such as A-MemGuard~\citep{amemguard} cross-check reasoning paths at run time. Storage-time provenance and policy-aware retention are inherited; the framework proposes something stronger and narrower: memory ownership is a cage on the transition relation, so that memory pollution is a reachability property that can be stated and, on an abstraction, checked, rather than a policy violation detected after the fact. This is argued, not formalised; the detectors are the natural comparison group.

\paragraph{Preregistration and falsifiability.}
Popper's demarcation~\citep{popper1959}, the preregistration movement~\citep{nosek2018}, registered reports~\citep{chambers2013}, and test-driven development's rule that the test precedes the code~\citep{beck2003}. The priority requirement is the same in each; the framework imports it unchanged into the design of an adjudicator, observes that the Meta-Cage is what makes it enforceable rather than aspirational, and in \S\ref{sec:experiments} applies it to itself.

\paragraph{Formal methods.}
Lean~4~\citep{lean4}, TLA$^{+}$~\citep{lamport2002}, and refinement mappings~\citep{abadi1991}. The posture is inherited: the facts the paper leans on are stated in checked form, and refinement is the correct form of the claim that a runtime implements caging. Nothing is claimed to be new here; by the framework's own lights, a statement about it that could be checked and has not been is a draft.
